\scanpage{64}

\mspage{5}

\uline{démonstration} : 1) on remarque d'abord que la conclusion ne dépend qu'apparamment du système de générateurs choisis : \unsure{o}

\[
\begin{aligned}
& f_1..f_n\ \text{régulier}\Longleftrightarrow f_1..f_n\ \text{quasi-regulier}\Longleftrightarrow\\
& A/I[T_1..T_n]\otimes_{A/I}M/I\longrightarrow\operatorname{gr}_I(A)\otimes_{\operatorname{gr}_I^0(A)}\operatorname{gr}_I^0(M)\longrightarrow\operatorname{gr}_I(M)\quad\text{bijectif}\\
\Longleftrightarrow{}&\operatorname{gr}_I(A)\otimes_{A/I}M/I\longrightarrow\operatorname{gr}_I(M)\quad\text{bijectif}\qquad\text{et}\\
& A/I[T_1..T_n]\otimes A/\operatorname{Ann}(M/I)\longrightarrow\operatorname{gr}_I(A)\otimes A/\operatorname{Ann}(M/I)\quad\text{bijectif.}\\
\Longleftrightarrow{}&\operatorname{gr}_I(A)\otimes M/I\xrightarrow{\sim}\operatorname{gr}_I(M),\quad\text{et \uline{modulo $\operatorname{Ann}(M/I)$}},\quad I/I^2\ \text{libre de rg }n\text{ et}\\
&\operatorname{Sym}_{A/I}(I/I^2)\xrightarrow{\sim}\operatorname{gr}_I(A).
\end{aligned}
\]

2) supposons qu'il existe $f$ dans $I\setminus mI$ qui soit $M$-régulier. on sait que $\operatorname{Ext}^i_A(A/I,M)$ \deleted{signifie} $=0$ pour $i<n$ signifie qu'il existe une $n$-suite $M$-régulière contenue dans $I$. On aura donc $\operatorname{Ext}^i_{A/f}(A/I,M/fM)=0$ pour $i\unsure{<|\leq}n-1$ et $I/f$ est engendré par $n-1$ générateurs. On conclut par récurrence sur $n$

3) supposons enfin qu'aucun élément de $I\setminus mI$ ne soit $M$-régulier. On a $I\subset mI\cup\bigcup_{p\in\operatorname{Ass}(M)}p$, donc (lemme plus bas) $\exists p\in\operatorname{Ass}M:I\subset p$. et $\operatorname{Hom}(A/I,M)\ne0$ donc $n=0$, $I=0$ et l'assertion est vraie

On a utilisé le

\uline{lemme} (Nagata !) : Soit $a$ un idéal d'un anneau $A$, $(b_i)_{0<i\leq n}$ des idéaux, $(p_i)_{0<i\leq m}$ des idéaux premiers. Si $a\subset\bigcup_i b_i\cup\bigcup_i p_i$, soit $a\subset\bigcup b_i$, soit $\exists i:a\subset p_i$.

% broken underline below the first word
\uline{démonstration} : $\alpha$) reduction à $m=1$ [regarder les $p$ comme $b$, sauf 1 et récurrer] et à ce que $p\not\supset b_i$. $\beta$) Si $a\not\subset\bigcup b_i$, soit $x\in a\setminus\bigcup b_i$ donc $x\in p$. Si $y\in a\cap\bigcap b_i$, $x+y\in a\setminus\bigcup b_i$ donc $x+y\in p$ et $y\in p$ : $a\cap\bigcap b_i\subset p$ et comme $b_i\not\subset p$, $a\subset p$.

\deleted{[NB : un dévissage analogue permet je crois de montrer (iv)$\Rightarrow$(iii) sans utiliser l'hypothèse \unsure{conjointe à|adjointe à} (iv)]}

$0_{\overline{\underline{\mathrm{IV}}}}$ 16.1.9.\qquad On peut y supprimer l'hypothèse noethérienne, en référant à Bourb. Alg Comm ch V § 1 n° 1 Th 1 ($E_{\mathrm{III}}$)

$0_{\overline{\underline{\mathrm{IV}}}}$ 19.2\qquad Il manque la proposition suivante, dont tu démontres un cas particulier en 19.7.1.1 (et un peu plus, car tu prouves là $B$ strictement formellement projectif)

\uline{prop} : Soit $A\longrightarrow B$ un homomorphisme \uline{local} d'\uline{anneaux locaux noethériens}, \uline{$M$} un \uline{$B$-module de type fini} ; on munit $A$ et $M$ des \uline{topologies $m$ et $n$-adiques} ($m=R(A),n=R(B)$. Les conditions suivantes sont \uline{équivalentes} :
% archive sheet number: 63
